% Source: https://artofproblemsolving.com/community/c3h1566519_geometryprovements
% Stable identity: aops-topic-1566519-q01
Let the common base be $BC$ and the common area $K$. The height to $BC$ is $h=\dfrac{2K}{BC}$ for all these triangles, so the third vertex $A$ lies on a line $\ell$ parallel to $BC$ at distance $h$ (on a fixed side of $BC$; the other side is symmetric). The perimeter is $BC+AB+AC$, so we must minimise $AB+AC$ for $A$ on $\ell$.

Reflect $B$ in $\ell$ to $B'$. For every $A$ on $\ell$, $AB=AB'$, so by the triangle inequality
\[
AB+AC=AB'+AC\ge B'C,
\]
with equality exactly when $A$ lies on segment $B'C$, i.e.\ at the point $A_0$ where $B'C$ crosses $\ell$.

\begin{asy}
usepackage("amsmath");
size(6.5cm);
defaultpen(fontsize(10pt));
pair B=(0,0),C=(6,0),Bp=(0,6);
pair A0=(3,3),A=(4.6,3);
draw((-1,3)--(7.5,3),gray);label("$\ell$",(7.5,3),E,gray);
draw(B--C,linewidth(0.9));
draw(B--A--C,blue+linewidth(0.9));
draw(B--A0--C,linewidth(0.6)+deepgreen);
draw(B--Bp,dotted);
draw(Bp--A,blue+dashed);draw(Bp--A0,red+dashed);
dot(A);dot(A0);dot(Bp);
label("$B$",B,SW);label("$C$",C,SE);label("$B'$",Bp,NW);
label("$A$",A,NE);label("$A_0$",A0,(0,-1.6));
label("$h$",(-0.25,1.5),W);label("$h$",(-0.25,4.5),W);
draw((-0.18,1.5)--(0.18,1.5));draw((-0.18,4.5)--(0.18,4.5));
label("\mbox{}",(min(currentpicture,user=true)+max(currentpicture,user=true))/2,black); // keeps the text after the figure black when LaTeX typesets the labels inline
\end{asy}

$BB'$ is perpendicular to $BC$ with $BB'=2h$, and $\ell$ is halfway between $B$ and $B'$. So $A_0$, being on $\ell$, is the midpoint of $B'C$, and its distance from line $BB'$ is half that of $C$, namely $\dfrac{BC}{2}$. Thus $A_0$ lies on the perpendicular bisector of $BC$, i.e.\ $A_0B=A_0C$.

Hence the minimum of $AB+AC$, and so of the perimeter, is attained exactly by the isosceles triangle $A_0BC$. $\blacksquare$
